% --- Archon physics formalization source begin ---
% archon:physics
% archon:covers PhyXMiniProblems/problem_phyx_mini_0935.lean
% archon:source-report reports/phyx_mini/problem_phyx_mini_0935.source.json

\chapter{Physics problem phyx\_mini\_0935}
\label{ch:PhyXMiniProblems_problem_phyx_mini_0935}

\paragraph{Problem source.}
\#\# Physical scenario

The magnetic field between the poles of the electromagnet in figure is uniform at any time, but its magnitude is increasing at the rate of 0.020 T/s. The area of the conducting loop in the field is \$120 \textbackslash{}, cm\textasciicircum{}2\$, and the total circuit resistance, including the meter, is \$5.0 \textbackslash{}, \textbackslash{}Omega\$.

\#\# Question

Find the induced emf and the induced current in the circuit.

\#\# Answer choices

- A: A: 0.45mV

- B: B: 0.24mV

- C: C: 0.54mV

- D: D: 0.23mV

\#\# Figure caption (auxiliary; use the image as primary evidence)

The image depicts a magnetic circuit setup. Here's a detailed description:

1. **Magnetic Field:**
   - Two gray cylindrical magnets are shown with the top labeled "S" (South) and the bottom labeled "N" (North). 
   - Blue arrows represent the magnetic field lines traveling from the South to the North pole.

2. **Loop and Area:**
   - A circular loop is placed between the magnets.
   - A vector labeled "A" points upwards inside the area of the loop, indicating the area vector of the loop.
   - The area of the loop is given as \textbackslash{}( A = 120 \textbackslash{}, \textbackslash{}text\{cm\}\textasciicircum{}2 = 0.012 \textbackslash{}, \textbackslash{}text\{m\}\textasciicircum{}2 \textbackslash{}).

3. **Current and Induction:**
   - The loop is connected to a wire, with points labeled "a" and "b".
   - A purple arrow labeled "I" indicates the direction of induced current, which is counterclockwise when viewed from above in this setup.

4. **Rate of Change of Magnetic Field:**
   - A formula is shown, \textbackslash{}(\textbackslash{}frac\{dB\}\{dt\} = 0.020 \textbackslash{}, \textbackslash{}text\{T/s\}\textbackslash{}), representing the rate of change of the magnetic field.

5. **Resistance:**
   - The total resistance in the circuit, including the meter, is given as \textbackslash{}(5.0 \textbackslash{}, \textbackslash{}Omega\textbackslash{}).

6. **Meter:**
   - A meter, possibly an ammeter or voltmeter, is connected to the circuit, with its needle slightly deflected.

This setup illustrates electromagnetic induction as described by Faraday's Law, where a changing magnetic field induces an electromotive force (EMF) and hence a current in the circuit.

\#\# Dataset metadata

Electromagnetic Induction; Physical Model Grounding Reasoning, Multi-Formula Reasoning

\paragraph{Recorded answer/context.}
B: B: 0.24mV

\paragraph{Figure/image path.}
/root/proposal\_for\_physic/science-mango/phyx\_mini\_run/phyx\_data/test\_image/935.png

\paragraph{Formalization target.}
create a compiling Lean file with sorry bodies at `PhyXMiniProblems/problem\_phyx\_mini\_0935.lean`. The Lean declarations must preserve the physical quantities, dimensions or dimensional roles, figure labels, governing-law hypotheses, and final relation expressed by this problem.
Use Mathlib/Physlib names found through LeanExplore where available. If a physics API is missing, introduce faithful local abstractions rather than scalar placeholder aliases.

\begin{theorem}[Physics formalization target]
\label{thm:physics:phyx_mini_0935:target}
\lean{PhyXMiniProblems.ProblemPhyXMini0935.inducedEmf_and_current_for_supplied_loop}
\uses{def:physics:phyx-mini-0935:phyxminiproblems-problemphyxmini0935-electromotiveforcemagnitudeinmillivolts, def:physics:phyx-mini-0935:phyxminiproblems-problemphyxmini0935-electriccurrentmagnitudeinmilliamperes, def:physics:phyx-mini-0935:phyxminiproblems-problemphyxmini0935-inductionloopsetup, def:physics:phyx-mini-0935:phyxminiproblems-problemphyxmini0935-matchesprimaryinductionfigure, def:physics:phyx-mini-0935:phyxminiproblems-problemphyxmini0935-matchesinductionproblemdescription, def:physics:phyx-mini-0935:phyxminiproblems-problemphyxmini0935-hasphysicalinductionparameters, def:physics:phyx-mini-0935:phyxminiproblems-problemphyxmini0935-satisfieselectromagneticinductionlaws, def:physics:phyx-mini-0935:phyxminiproblems-problemphyxmini0935-recordeddatasetanswer, def:physics:phyx-mini-0935:phyxminiproblems-problemphyxmini0935-answermatchesinducedemf, def:physics:phyx-mini-0935:phyxminiproblems-problemphyxmini0935-inducedcurrentisclockwiseviewedfromabove}
The assigned autoformalize agent should translate this physics problem into Lean declarations in the covered file, with theorem and lemma proof bodies written as `by sorry`.
\end{theorem}
\begin{proof}
This is an autoformalization task, not a proof task. Produce statements that can later be proved by the physics prover without weakening the physical model.
\end{proof}

% --- Archon declaration topology begin ---
\section{Declaration topology}
\begin{definition}[electrical Resistance Dimension]
  \label{def:physics:phyx-mini-0935:phyxminiproblems-problemphyxmini0935-electricalresistancedimension}
  \lean{PhyXMiniProblems.ProblemPhyXMini0935.electricalResistanceDimension}
  Physical dimension M L² T⁻¹ C⁻² of electrical resistance
\end{definition}
\begin{proof}
  This is the declared model definition of electrical Resistance Dimension.
\end{proof}
\begin{definition}[magnetic Flux Density Dimension]
  \label{def:physics:phyx-mini-0935:phyxminiproblems-problemphyxmini0935-magneticfluxdensitydimension}
  \lean{PhyXMiniProblems.ProblemPhyXMini0935.magneticFluxDensityDimension}
  Physical dimension M T⁻¹ C⁻¹ of magnetic flux density (tesla)
\end{definition}
\begin{proof}
  This is the declared model definition of magnetic Flux Density Dimension.
\end{proof}
\begin{definition}[magnetic Flux Density Rate Dimension]
  \label{def:physics:phyx-mini-0935:phyxminiproblems-problemphyxmini0935-magneticfluxdensityratedimension}
  \lean{PhyXMiniProblems.ProblemPhyXMini0935.magneticFluxDensityRateDimension}
  \uses{def:physics:phyx-mini-0935:phyxminiproblems-problemphyxmini0935-magneticfluxdensitydimension}
  Physical dimension M T⁻² C⁻¹ of magnetic-flux-density rate
\end{definition}
\begin{proof}
  This is the declared model definition of magnetic Flux Density Rate Dimension.
\end{proof}
\begin{definition}[magnetic Flux Rate Dimension]
  \label{def:physics:phyx-mini-0935:phyxminiproblems-problemphyxmini0935-magneticfluxratedimension}
  \lean{PhyXMiniProblems.ProblemPhyXMini0935.magneticFluxRateDimension}
  \uses{def:physics:phyx-mini-0935:phyxminiproblems-problemphyxmini0935-magneticfluxdensityratedimension}
  Physical dimension M L² T⁻² C⁻¹ of magnetic-flux rate
\end{definition}
\begin{proof}
  This is the declared model definition of magnetic Flux Rate Dimension.
\end{proof}
\begin{definition}[electromotive Force Dimension]
  \label{def:physics:phyx-mini-0935:phyxminiproblems-problemphyxmini0935-electromotiveforcedimension}
  \lean{PhyXMiniProblems.ProblemPhyXMini0935.electromotiveForceDimension}
  Physical dimension M L² T⁻² C⁻¹ of electromotive force
\end{definition}
\begin{proof}
  This is the declared model definition of electromotive Force Dimension.
\end{proof}
\begin{definition}[electric Current Dimension]
  \label{def:physics:phyx-mini-0935:phyxminiproblems-problemphyxmini0935-electriccurrentdimension}
  \lean{PhyXMiniProblems.ProblemPhyXMini0935.electricCurrentDimension}
  Physical dimension C T⁻¹ of electric current
\end{definition}
\begin{proof}
  This is the declared model definition of electric Current Dimension.
\end{proof}
\begin{definition}[Area Magnitude]
  \label{def:physics:phyx-mini-0935:phyxminiproblems-problemphyxmini0935-areamagnitude}
  \lean{PhyXMiniProblems.ProblemPhyXMini0935.AreaMagnitude}
  A nonnegative, unit-independent loop area from Physlib's area API
\end{definition}
\begin{proof}
  This is the declared model definition of Area Magnitude.
\end{proof}
\begin{definition}[Electrical Resistance Magnitude]
  \label{def:physics:phyx-mini-0935:phyxminiproblems-problemphyxmini0935-electricalresistancemagnitude}
  \lean{PhyXMiniProblems.ProblemPhyXMini0935.ElectricalResistanceMagnitude}
  \uses{def:physics:phyx-mini-0935:phyxminiproblems-problemphyxmini0935-electricalresistancedimension}
  A nonnegative, unit-independent electrical resistance
\end{definition}
\begin{proof}
  This is the declared model definition of Electrical Resistance Magnitude.
\end{proof}
\begin{definition}[Magnetic Flux Density Magnitude]
  \label{def:physics:phyx-mini-0935:phyxminiproblems-problemphyxmini0935-magneticfluxdensitymagnitude}
  \lean{PhyXMiniProblems.ProblemPhyXMini0935.MagneticFluxDensityMagnitude}
  \uses{def:physics:phyx-mini-0935:phyxminiproblems-problemphyxmini0935-magneticfluxdensitydimension}
  A nonnegative, unit-independent magnetic-flux-density magnitude
\end{definition}
\begin{proof}
  This is the declared model definition of Magnetic Flux Density Magnitude.
\end{proof}
\begin{definition}[Magnetic Flux Density Rate Magnitude]
  \label{def:physics:phyx-mini-0935:phyxminiproblems-problemphyxmini0935-magneticfluxdensityratemagnitude}
  \lean{PhyXMiniProblems.ProblemPhyXMini0935.MagneticFluxDensityRateMagnitude}
  \uses{def:physics:phyx-mini-0935:phyxminiproblems-problemphyxmini0935-magneticfluxdensityratedimension}
  A nonnegative, unit-independent magnetic-flux-density rate
\end{definition}
\begin{proof}
  This is the declared model definition of Magnetic Flux Density Rate Magnitude.
\end{proof}
\begin{definition}[Signed Magnetic Flux Rate]
  \label{def:physics:phyx-mini-0935:phyxminiproblems-problemphyxmini0935-signedmagneticfluxrate}
  \lean{PhyXMiniProblems.ProblemPhyXMini0935.SignedMagneticFluxRate}
  \uses{def:physics:phyx-mini-0935:phyxminiproblems-problemphyxmini0935-magneticfluxratedimension}
  A signed, unit-independent magnetic-flux rate
\end{definition}
\begin{proof}
  This is the declared model definition of Signed Magnetic Flux Rate.
\end{proof}
\begin{definition}[Electromotive Force]
  \label{def:physics:phyx-mini-0935:phyxminiproblems-problemphyxmini0935-electromotiveforce}
  \lean{PhyXMiniProblems.ProblemPhyXMini0935.ElectromotiveForce}
  \uses{def:physics:phyx-mini-0935:phyxminiproblems-problemphyxmini0935-electromotiveforcedimension}
  A signed, unit-independent electromotive force
\end{definition}
\begin{proof}
  This is the declared model definition of Electromotive Force.
\end{proof}
\begin{definition}[Electric Current]
  \label{def:physics:phyx-mini-0935:phyxminiproblems-problemphyxmini0935-electriccurrent}
  \lean{PhyXMiniProblems.ProblemPhyXMini0935.ElectricCurrent}
  \uses{def:physics:phyx-mini-0935:phyxminiproblems-problemphyxmini0935-electriccurrentdimension}
  A signed, unit-independent electric current
\end{definition}
\begin{proof}
  This is the declared model definition of Electric Current.
\end{proof}
\begin{definition}[area Readout]
  \label{def:physics:phyx-mini-0935:phyxminiproblems-problemphyxmini0935-areareadout}
  \lean{PhyXMiniProblems.ProblemPhyXMini0935.areaReadout}
  \uses{def:physics:phyx-mini-0935:phyxminiproblems-problemphyxmini0935-areamagnitude}
  Read physical area in the coherent area unit selected by units
\end{definition}
\begin{proof}
  This is the declared model definition of area Readout.
\end{proof}
\begin{definition}[resistance Readout]
  \label{def:physics:phyx-mini-0935:phyxminiproblems-problemphyxmini0935-resistancereadout}
  \lean{PhyXMiniProblems.ProblemPhyXMini0935.resistanceReadout}
  \uses{def:physics:phyx-mini-0935:phyxminiproblems-problemphyxmini0935-electricalresistancemagnitude}
  Read electrical resistance in the coherent unit selected by units
\end{definition}
\begin{proof}
  This is the declared model definition of resistance Readout.
\end{proof}
\begin{definition}[magnetic Flux Density Readout]
  \label{def:physics:phyx-mini-0935:phyxminiproblems-problemphyxmini0935-magneticfluxdensityreadout}
  \lean{PhyXMiniProblems.ProblemPhyXMini0935.magneticFluxDensityReadout}
  \uses{def:physics:phyx-mini-0935:phyxminiproblems-problemphyxmini0935-magneticfluxdensitymagnitude}
  Read magnetic-flux-density magnitude in the selected coherent unit
\end{definition}
\begin{proof}
  This is the declared model definition of magnetic Flux Density Readout.
\end{proof}
\begin{definition}[magnetic Flux Density Rate Readout]
  \label{def:physics:phyx-mini-0935:phyxminiproblems-problemphyxmini0935-magneticfluxdensityratereadout}
  \lean{PhyXMiniProblems.ProblemPhyXMini0935.magneticFluxDensityRateReadout}
  \uses{def:physics:phyx-mini-0935:phyxminiproblems-problemphyxmini0935-magneticfluxdensityratemagnitude}
  Read magnetic-flux-density rate in the selected coherent unit
\end{definition}
\begin{proof}
  This is the declared model definition of magnetic Flux Density Rate Readout.
\end{proof}
\begin{definition}[signed Magnetic Flux Rate Readout]
  \label{def:physics:phyx-mini-0935:phyxminiproblems-problemphyxmini0935-signedmagneticfluxratereadout}
  \lean{PhyXMiniProblems.ProblemPhyXMini0935.signedMagneticFluxRateReadout}
  \uses{def:physics:phyx-mini-0935:phyxminiproblems-problemphyxmini0935-signedmagneticfluxrate}
  Read signed magnetic-flux rate in the selected coherent unit
\end{definition}
\begin{proof}
  This is the declared model definition of signed Magnetic Flux Rate Readout.
\end{proof}
\begin{definition}[electromotive Force Readout]
  \label{def:physics:phyx-mini-0935:phyxminiproblems-problemphyxmini0935-electromotiveforcereadout}
  \lean{PhyXMiniProblems.ProblemPhyXMini0935.electromotiveForceReadout}
  \uses{def:physics:phyx-mini-0935:phyxminiproblems-problemphyxmini0935-electromotiveforce}
  Read signed electromotive force in the selected coherent unit
\end{definition}
\begin{proof}
  This is the declared model definition of electromotive Force Readout.
\end{proof}
\begin{definition}[electric Current Readout]
  \label{def:physics:phyx-mini-0935:phyxminiproblems-problemphyxmini0935-electriccurrentreadout}
  \lean{PhyXMiniProblems.ProblemPhyXMini0935.electricCurrentReadout}
  \uses{def:physics:phyx-mini-0935:phyxminiproblems-problemphyxmini0935-electriccurrent}
  Read signed electric current in the selected coherent unit
\end{definition}
\begin{proof}
  This is the declared model definition of electric Current Readout.
\end{proof}
\begin{definition}[area In Square Meters]
  \label{def:physics:phyx-mini-0935:phyxminiproblems-problemphyxmini0935-areainsquaremeters}
  \lean{PhyXMiniProblems.ProblemPhyXMini0935.areaInSquareMeters}
  \uses{def:physics:phyx-mini-0935:phyxminiproblems-problemphyxmini0935-areamagnitude, def:physics:phyx-mini-0935:phyxminiproblems-problemphyxmini0935-areareadout}
  Coherent-SI square-metre readout of loop area
\end{definition}
\begin{proof}
  This is the declared model definition of area In Square Meters.
\end{proof}
\begin{definition}[area In Square Centimeters]
  \label{def:physics:phyx-mini-0935:phyxminiproblems-problemphyxmini0935-areainsquarecentimeters}
  \lean{PhyXMiniProblems.ProblemPhyXMini0935.areaInSquareCentimeters}
  \uses{def:physics:phyx-mini-0935:phyxminiproblems-problemphyxmini0935-areamagnitude, def:physics:phyx-mini-0935:phyxminiproblems-problemphyxmini0935-areainsquaremeters}
  Square-centimetre readout used by the printed area label
\end{definition}
\begin{proof}
  This is the declared model definition of area In Square Centimeters.
\end{proof}
\begin{definition}[resistance In Ohms]
  \label{def:physics:phyx-mini-0935:phyxminiproblems-problemphyxmini0935-resistanceinohms}
  \lean{PhyXMiniProblems.ProblemPhyXMini0935.resistanceInOhms}
  \uses{def:physics:phyx-mini-0935:phyxminiproblems-problemphyxmini0935-electricalresistancemagnitude, def:physics:phyx-mini-0935:phyxminiproblems-problemphyxmini0935-resistancereadout}
  Coherent-SI ohm readout of total circuit resistance
\end{definition}
\begin{proof}
  This is the declared model definition of resistance In Ohms.
\end{proof}
\begin{definition}[magnetic Flux Density In Teslas]
  \label{def:physics:phyx-mini-0935:phyxminiproblems-problemphyxmini0935-magneticfluxdensityinteslas}
  \lean{PhyXMiniProblems.ProblemPhyXMini0935.magneticFluxDensityInTeslas}
  \uses{def:physics:phyx-mini-0935:phyxminiproblems-problemphyxmini0935-magneticfluxdensitymagnitude, def:physics:phyx-mini-0935:phyxminiproblems-problemphyxmini0935-magneticfluxdensityreadout}
  Coherent-SI tesla readout of magnetic-flux-density magnitude
\end{definition}
\begin{proof}
  This is the declared model definition of magnetic Flux Density In Teslas.
\end{proof}
\begin{definition}[magnetic Flux Density Rate In Teslas Per Second]
  \label{def:physics:phyx-mini-0935:phyxminiproblems-problemphyxmini0935-magneticfluxdensityrateinteslaspersecond}
  \lean{PhyXMiniProblems.ProblemPhyXMini0935.magneticFluxDensityRateInTeslasPerSecond}
  \uses{def:physics:phyx-mini-0935:phyxminiproblems-problemphyxmini0935-magneticfluxdensityratemagnitude, def:physics:phyx-mini-0935:phyxminiproblems-problemphyxmini0935-magneticfluxdensityratereadout}
  Coherent-SI tesla-per-second readout of field-magnitude rate
\end{definition}
\begin{proof}
  This is the declared model definition of magnetic Flux Density Rate In Teslas Per Second.
\end{proof}
\begin{definition}[signed Magnetic Flux Rate In Webers Per Second]
  \label{def:physics:phyx-mini-0935:phyxminiproblems-problemphyxmini0935-signedmagneticfluxrateinweberspersecond}
  \lean{PhyXMiniProblems.ProblemPhyXMini0935.signedMagneticFluxRateInWebersPerSecond}
  \uses{def:physics:phyx-mini-0935:phyxminiproblems-problemphyxmini0935-signedmagneticfluxrate, def:physics:phyx-mini-0935:phyxminiproblems-problemphyxmini0935-signedmagneticfluxratereadout}
  Coherent-SI weber-per-second readout of signed magnetic-flux rate
\end{definition}
\begin{proof}
  This is the declared model definition of signed Magnetic Flux Rate In Webers Per Second.
\end{proof}
\begin{definition}[electromotive Force In Volts]
  \label{def:physics:phyx-mini-0935:phyxminiproblems-problemphyxmini0935-electromotiveforceinvolts}
  \lean{PhyXMiniProblems.ProblemPhyXMini0935.electromotiveForceInVolts}
  \uses{def:physics:phyx-mini-0935:phyxminiproblems-problemphyxmini0935-electromotiveforce, def:physics:phyx-mini-0935:phyxminiproblems-problemphyxmini0935-electromotiveforcereadout}
  Coherent-SI volt readout of signed electromotive force
\end{definition}
\begin{proof}
  This is the declared model definition of electromotive Force In Volts.
\end{proof}
\begin{definition}[electromotive Force In Millivolts]
  \label{def:physics:phyx-mini-0935:phyxminiproblems-problemphyxmini0935-electromotiveforceinmillivolts}
  \lean{PhyXMiniProblems.ProblemPhyXMini0935.electromotiveForceInMillivolts}
  \uses{def:physics:phyx-mini-0935:phyxminiproblems-problemphyxmini0935-electromotiveforce, def:physics:phyx-mini-0935:phyxminiproblems-problemphyxmini0935-electromotiveforceinvolts}
  Signed millivolt readout relative to the drawn current traversal
\end{definition}
\begin{proof}
  This is the declared model definition of electromotive Force In Millivolts.
\end{proof}
\begin{definition}[electromotive Force Magnitude In Millivolts]
  \label{def:physics:phyx-mini-0935:phyxminiproblems-problemphyxmini0935-electromotiveforcemagnitudeinmillivolts}
  \lean{PhyXMiniProblems.ProblemPhyXMini0935.electromotiveForceMagnitudeInMillivolts}
  \uses{def:physics:phyx-mini-0935:phyxminiproblems-problemphyxmini0935-electromotiveforce, def:physics:phyx-mini-0935:phyxminiproblems-problemphyxmini0935-electromotiveforceinmillivolts}
  The unsigned emf magnitude requested by the positive answer choices
\end{definition}
\begin{proof}
  This is the declared model definition of electromotive Force Magnitude In Millivolts.
\end{proof}
\begin{definition}[electric Current In Amperes]
  \label{def:physics:phyx-mini-0935:phyxminiproblems-problemphyxmini0935-electriccurrentinamperes}
  \lean{PhyXMiniProblems.ProblemPhyXMini0935.electricCurrentInAmperes}
  \uses{def:physics:phyx-mini-0935:phyxminiproblems-problemphyxmini0935-electriccurrent, def:physics:phyx-mini-0935:phyxminiproblems-problemphyxmini0935-electriccurrentreadout}
  Coherent-SI ampere readout, signed relative to the drawn arrow
\end{definition}
\begin{proof}
  This is the declared model definition of electric Current In Amperes.
\end{proof}
\begin{definition}[electric Current In Milliamperes]
  \label{def:physics:phyx-mini-0935:phyxminiproblems-problemphyxmini0935-electriccurrentinmilliamperes}
  \lean{PhyXMiniProblems.ProblemPhyXMini0935.electricCurrentInMilliamperes}
  \uses{def:physics:phyx-mini-0935:phyxminiproblems-problemphyxmini0935-electriccurrent, def:physics:phyx-mini-0935:phyxminiproblems-problemphyxmini0935-electriccurrentinamperes}
  Milliampere readout, signed relative to the drawn arrow
\end{definition}
\begin{proof}
  This is the declared model definition of electric Current In Milliamperes.
\end{proof}
\begin{definition}[electric Current Magnitude In Milliamperes]
  \label{def:physics:phyx-mini-0935:phyxminiproblems-problemphyxmini0935-electriccurrentmagnitudeinmilliamperes}
  \lean{PhyXMiniProblems.ProblemPhyXMini0935.electricCurrentMagnitudeInMilliamperes}
  \uses{def:physics:phyx-mini-0935:phyxminiproblems-problemphyxmini0935-electriccurrent, def:physics:phyx-mini-0935:phyxminiproblems-problemphyxmini0935-electriccurrentinmilliamperes}
  The unsigned current magnitude requested by the problem
\end{definition}
\begin{proof}
  This is the declared model definition of electric Current Magnitude In Milliamperes.
\end{proof}
\begin{definition}[Magnetic Pole]
  \label{def:physics:phyx-mini-0935:phyxminiproblems-problemphyxmini0935-magneticpole}
  \lean{PhyXMiniProblems.ProblemPhyXMini0935.MagneticPole}
  Labels printed on the two electromagnet pole pieces
\end{definition}
\begin{proof}
  This is the declared model definition of Magnetic Pole.
\end{proof}
\begin{definition}[Pole Position]
  \label{def:physics:phyx-mini-0935:phyxminiproblems-problemphyxmini0935-poleposition}
  \lean{PhyXMiniProblems.ProblemPhyXMini0935.PolePosition}
  Vertical locations of the two pole pieces in the primary image
\end{definition}
\begin{proof}
  This is the declared model definition of Pole Position.
\end{proof}
\begin{definition}[Loop Point]
  \label{def:physics:phyx-mini-0935:phyxminiproblems-problemphyxmini0935-looppoint}
  \lean{PhyXMiniProblems.ProblemPhyXMini0935.LoopPoint}
  The two lettered points where the external circuit meets the loop
\end{definition}
\begin{proof}
  This is the declared model definition of Loop Point.
\end{proof}
\begin{definition}[Vertical Direction]
  \label{def:physics:phyx-mini-0935:phyxminiproblems-problemphyxmini0935-verticaldirection}
  \lean{PhyXMiniProblems.ProblemPhyXMini0935.VerticalDirection}
  Vertical directions available to field and area-vector arrows
\end{definition}
\begin{proof}
  This is the declared model definition of Vertical Direction.
\end{proof}
\begin{definition}[Loop Traversal Direction]
  \label{def:physics:phyx-mini-0935:phyxminiproblems-problemphyxmini0935-looptraversaldirection}
  \lean{PhyXMiniProblems.ProblemPhyXMini0935.LoopTraversalDirection}
  Traversal directions at the visible front portion of the loop
\end{definition}
\begin{proof}
  This is the declared model definition of Loop Traversal Direction.
\end{proof}
\begin{definition}[Loop Sense From Above]
  \label{def:physics:phyx-mini-0935:phyxminiproblems-problemphyxmini0935-loopsensefromabove}
  \lean{PhyXMiniProblems.ProblemPhyXMini0935.LoopSenseFromAbove}
  Sense in which a horizontal loop is traversed when viewed from above
\end{definition}
\begin{proof}
  This is the declared model definition of Loop Sense From Above.
\end{proof}
\begin{definition}[Axial Alignment]
  \label{def:physics:phyx-mini-0935:phyxminiproblems-problemphyxmini0935-axialalignment}
  \lean{PhyXMiniProblems.ProblemPhyXMini0935.AxialAlignment}
  Parallel or antiparallel alignment of two axial directions
\end{definition}
\begin{proof}
  This is the declared model definition of Axial Alignment.
\end{proof}
\begin{definition}[Loop Topology]
  \label{def:physics:phyx-mini-0935:phyxminiproblems-problemphyxmini0935-looptopology}
  \lean{PhyXMiniProblems.ProblemPhyXMini0935.LoopTopology}
  Topology and conductor count of the depicted loop
\end{definition}
\begin{proof}
  This is the declared model definition of Loop Topology.
\end{proof}
\begin{definition}[Loop Motion Model]
  \label{def:physics:phyx-mini-0935:phyxminiproblems-problemphyxmini0935-loopmotionmodel}
  \lean{PhyXMiniProblems.ProblemPhyXMini0935.LoopMotionModel}
  Mechanical state relevant to the motional-emf contribution
\end{definition}
\begin{proof}
  This is the declared model definition of Loop Motion Model.
\end{proof}
\begin{definition}[Meter Kind]
  \label{def:physics:phyx-mini-0935:phyxminiproblems-problemphyxmini0935-meterkind}
  \lean{PhyXMiniProblems.ProblemPhyXMini0935.MeterKind}
  Kind of meter identifiable from the primary image
\end{definition}
\begin{proof}
  This is the declared model definition of Meter Kind.
\end{proof}
\begin{definition}[expected Pole At]
  \label{def:physics:phyx-mini-0935:phyxminiproblems-problemphyxmini0935-expectedpoleat}
  \lean{PhyXMiniProblems.ProblemPhyXMini0935.expectedPoleAt}
  \uses{def:physics:phyx-mini-0935:phyxminiproblems-problemphyxmini0935-magneticpole, def:physics:phyx-mini-0935:phyxminiproblems-problemphyxmini0935-poleposition}
  Pole label expected at each location in image 935.png
\end{definition}
\begin{proof}
  This is the declared model definition of expected Pole At.
\end{proof}
\begin{definition}[right Hand Normal]
  \label{def:physics:phyx-mini-0935:phyxminiproblems-problemphyxmini0935-righthandnormal}
  \lean{PhyXMiniProblems.ProblemPhyXMini0935.rightHandNormal}
  \uses{def:physics:phyx-mini-0935:phyxminiproblems-problemphyxmini0935-verticaldirection, def:physics:phyx-mini-0935:phyxminiproblems-problemphyxmini0935-looptraversaldirection}
  Right-hand normal generated by a chosen positive loop traversal
\end{definition}
\begin{proof}
  This is the declared model definition of right Hand Normal.
\end{proof}
\begin{definition}[traversal Sense From Above]
  \label{def:physics:phyx-mini-0935:phyxminiproblems-problemphyxmini0935-traversalsensefromabove}
  \lean{PhyXMiniProblems.ProblemPhyXMini0935.traversalSenseFromAbove}
  \uses{def:physics:phyx-mini-0935:phyxminiproblems-problemphyxmini0935-looptraversaldirection, def:physics:phyx-mini-0935:phyxminiproblems-problemphyxmini0935-loopsensefromabove}
  View-from-above sense of the traversal shown at the loop's front
\end{definition}
\begin{proof}
  This is the declared model definition of traversal Sense From Above.
\end{proof}
\begin{definition}[axial Alignment]
  \label{def:physics:phyx-mini-0935:phyxminiproblems-problemphyxmini0935-axialalignment-value}
  \lean{PhyXMiniProblems.ProblemPhyXMini0935.axialAlignment}
  \uses{def:physics:phyx-mini-0935:phyxminiproblems-problemphyxmini0935-verticaldirection, def:physics:phyx-mini-0935:phyxminiproblems-problemphyxmini0935-axialalignment}
  Whether two vertical axial directions are parallel or antiparallel
\end{definition}
\begin{proof}
  This is the declared model definition of axial Alignment.
\end{proof}
\begin{definition}[cosine]
  \label{def:physics:phyx-mini-0935:phyxminiproblems-problemphyxmini0935-axialalignment-cosine}
  \lean{PhyXMiniProblems.ProblemPhyXMini0935.AxialAlignment.cosine}
  \uses{def:physics:phyx-mini-0935:phyxminiproblems-problemphyxmini0935-axialalignment}
  Cosine factor for the two axial alignments in this setup
\end{definition}
\begin{proof}
  This is the declared model definition of cosine.
\end{proof}
\begin{definition}[Electromagnetic Induction Figure]
  \label{def:physics:phyx-mini-0935:phyxminiproblems-problemphyxmini0935-electromagneticinductionfigure}
  \lean{PhyXMiniProblems.ProblemPhyXMini0935.ElectromagneticInductionFigure}
  \uses{def:physics:phyx-mini-0935:phyxminiproblems-problemphyxmini0935-magneticpole, def:physics:phyx-mini-0935:phyxminiproblems-problemphyxmini0935-poleposition, def:physics:phyx-mini-0935:phyxminiproblems-problemphyxmini0935-looppoint, def:physics:phyx-mini-0935:phyxminiproblems-problemphyxmini0935-verticaldirection, def:physics:phyx-mini-0935:phyxminiproblems-problemphyxmini0935-looptraversaldirection, def:physics:phyx-mini-0935:phyxminiproblems-problemphyxmini0935-meterkind}
  Literal geometric, presentation, and numerical content transcribed from the primary raster. The purple current arrow chooses a sign reference; it does not assume that the induced current actually follows the arrow
\end{definition}
\begin{proof}
  This is the declared model definition of Electromagnetic Induction Figure.
\end{proof}
\begin{definition}[Induction Loop Setup]
  \label{def:physics:phyx-mini-0935:phyxminiproblems-problemphyxmini0935-inductionloopsetup}
  \lean{PhyXMiniProblems.ProblemPhyXMini0935.InductionLoopSetup}
  \uses{def:physics:phyx-mini-0935:phyxminiproblems-problemphyxmini0935-areamagnitude, def:physics:phyx-mini-0935:phyxminiproblems-problemphyxmini0935-electricalresistancemagnitude, def:physics:phyx-mini-0935:phyxminiproblems-problemphyxmini0935-magneticfluxdensitymagnitude, def:physics:phyx-mini-0935:phyxminiproblems-problemphyxmini0935-magneticfluxdensityratemagnitude, def:physics:phyx-mini-0935:phyxminiproblems-problemphyxmini0935-signedmagneticfluxrate, def:physics:phyx-mini-0935:phyxminiproblems-problemphyxmini0935-electromotiveforce, def:physics:phyx-mini-0935:phyxminiproblems-problemphyxmini0935-electriccurrent, def:physics:phyx-mini-0935:phyxminiproblems-problemphyxmini0935-verticaldirection, def:physics:phyx-mini-0935:phyxminiproblems-problemphyxmini0935-looptraversaldirection, def:physics:phyx-mini-0935:phyxminiproblems-problemphyxmini0935-looptopology, def:physics:phyx-mini-0935:phyxminiproblems-problemphyxmini0935-loopmotionmodel, def:physics:phyx-mini-0935:phyxminiproblems-problemphyxmini0935-electromagneticinductionfigure}
  Independent physical objects and observables. The Physlib magnetic field, its dimensionful magnitude profile, induced emf, and induced current are not defined from numerical answers
\end{definition}
\begin{proof}
  This is the declared model definition of Induction Loop Setup.
\end{proof}
\begin{definition}[Matches Primary Induction Figure]
  \label{def:physics:phyx-mini-0935:phyxminiproblems-problemphyxmini0935-matchesprimaryinductionfigure}
  \lean{PhyXMiniProblems.ProblemPhyXMini0935.MatchesPrimaryInductionFigure}
  \uses{def:physics:phyx-mini-0935:phyxminiproblems-problemphyxmini0935-expectedpoleat, def:physics:phyx-mini-0935:phyxminiproblems-problemphyxmini0935-inductionloopsetup}
  Exact labels, arrows, components, and printed values in image 935.png
\end{definition}
\begin{proof}
  This is the declared model definition of Matches Primary Induction Figure.
\end{proof}
\begin{definition}[Matches Induction Problem Description]
  \label{def:physics:phyx-mini-0935:phyxminiproblems-problemphyxmini0935-matchesinductionproblemdescription}
  \lean{PhyXMiniProblems.ProblemPhyXMini0935.MatchesInductionProblemDescription}
  \uses{def:physics:phyx-mini-0935:phyxminiproblems-problemphyxmini0935-areainsquaremeters, def:physics:phyx-mini-0935:phyxminiproblems-problemphyxmini0935-areainsquarecentimeters, def:physics:phyx-mini-0935:phyxminiproblems-problemphyxmini0935-resistanceinohms, def:physics:phyx-mini-0935:phyxminiproblems-problemphyxmini0935-magneticfluxdensityinteslas, def:physics:phyx-mini-0935:phyxminiproblems-problemphyxmini0935-magneticfluxdensityrateinteslaspersecond, def:physics:phyx-mini-0935:phyxminiproblems-problemphyxmini0935-inductionloopsetup}
  Qualitative setup and calibration of independent physical quantities to the primary-image readouts. The Physlib field uses coordinates chosen so that its vector norm is read in teslas
\end{definition}
\begin{proof}
  This is the declared model definition of Matches Induction Problem Description.
\end{proof}
\begin{definition}[Has Physical Induction Parameters]
  \label{def:physics:phyx-mini-0935:phyxminiproblems-problemphyxmini0935-hasphysicalinductionparameters}
  \lean{PhyXMiniProblems.ProblemPhyXMini0935.HasPhysicalInductionParameters}
  \uses{def:physics:phyx-mini-0935:phyxminiproblems-problemphyxmini0935-areainsquaremeters, def:physics:phyx-mini-0935:phyxminiproblems-problemphyxmini0935-resistanceinohms, def:physics:phyx-mini-0935:phyxminiproblems-problemphyxmini0935-magneticfluxdensityrateinteslaspersecond, def:physics:phyx-mini-0935:phyxminiproblems-problemphyxmini0935-inductionloopsetup}
  Positivity conditions selecting the ordinary passive-circuit branch
\end{definition}
\begin{proof}
  This is the declared model definition of Has Physical Induction Parameters.
\end{proof}
\begin{definition}[Satisfies Electromagnetic Induction Laws]
  \label{def:physics:phyx-mini-0935:phyxminiproblems-problemphyxmini0935-satisfieselectromagneticinductionlaws}
  \lean{PhyXMiniProblems.ProblemPhyXMini0935.SatisfiesElectromagneticInductionLaws}
  \uses{def:physics:phyx-mini-0935:phyxminiproblems-problemphyxmini0935-areainsquaremeters, def:physics:phyx-mini-0935:phyxminiproblems-problemphyxmini0935-resistanceinohms, def:physics:phyx-mini-0935:phyxminiproblems-problemphyxmini0935-magneticfluxdensityinteslas, def:physics:phyx-mini-0935:phyxminiproblems-problemphyxmini0935-magneticfluxdensityrateinteslaspersecond, def:physics:phyx-mini-0935:phyxminiproblems-problemphyxmini0935-signedmagneticfluxrateinweberspersecond, def:physics:phyx-mini-0935:phyxminiproblems-problemphyxmini0935-electromotiveforceinvolts, def:physics:phyx-mini-0935:phyxminiproblems-problemphyxmini0935-electriccurrentinamperes, def:physics:phyx-mini-0935:phyxminiproblems-problemphyxmini0935-righthandnormal, def:physics:phyx-mini-0935:phyxminiproblems-problemphyxmini0935-axialalignment-value, def:physics:phyx-mini-0935:phyxminiproblems-problemphyxmini0935-inductionloopsetup}
  The governing induction and circuit laws. The signed flux rate uses the normal determined by the current sign reference, which is intentionally distinct from the independently displayed upward area vector. No requested emf or current value occurs in these laws
\end{definition}
\begin{proof}
  This is the declared model definition of Satisfies Electromagnetic Induction Laws.
\end{proof}
\begin{definition}[Answer Choice]
  \label{def:physics:phyx-mini-0935:phyxminiproblems-problemphyxmini0935-answerchoice}
  \lean{PhyXMiniProblems.ProblemPhyXMini0935.AnswerChoice}
  The four emf magnitudes, in millivolts, printed with the problem
\end{definition}
\begin{proof}
  This is the declared model definition of Answer Choice.
\end{proof}
\begin{definition}[displayed Emf In Millivolts]
  \label{def:physics:phyx-mini-0935:phyxminiproblems-problemphyxmini0935-answerchoice-displayedemfinmillivolts}
  \lean{PhyXMiniProblems.ProblemPhyXMini0935.AnswerChoice.displayedEmfInMillivolts}
  \uses{def:physics:phyx-mini-0935:phyxminiproblems-problemphyxmini0935-answerchoice}
  Emf magnitude in millivolts printed beside an answer choice
\end{definition}
\begin{proof}
  This is the declared model definition of displayed Emf In Millivolts.
\end{proof}
\begin{definition}[recorded Dataset Answer]
  \label{def:physics:phyx-mini-0935:phyxminiproblems-problemphyxmini0935-recordeddatasetanswer}
  \lean{PhyXMiniProblems.ProblemPhyXMini0935.recordedDatasetAnswer}
  \uses{def:physics:phyx-mini-0935:phyxminiproblems-problemphyxmini0935-answerchoice}
  Answer label recorded in the supplied dataset
\end{definition}
\begin{proof}
  This is the declared model definition of recorded Dataset Answer.
\end{proof}
\begin{definition}[Answer Matches Induced Emf]
  \label{def:physics:phyx-mini-0935:phyxminiproblems-problemphyxmini0935-answermatchesinducedemf}
  \lean{PhyXMiniProblems.ProblemPhyXMini0935.AnswerMatchesInducedEmf}
  \uses{def:physics:phyx-mini-0935:phyxminiproblems-problemphyxmini0935-electromotiveforcemagnitudeinmillivolts, def:physics:phyx-mini-0935:phyxminiproblems-problemphyxmini0935-inductionloopsetup, def:physics:phyx-mini-0935:phyxminiproblems-problemphyxmini0935-answerchoice}
  A choice matches the independently modeled induced-emf magnitude
\end{definition}
\begin{proof}
  This is the declared model definition of Answer Matches Induced Emf.
\end{proof}
\begin{definition}[Induced Current Runs Along Drawn Arrow]
  \label{def:physics:phyx-mini-0935:phyxminiproblems-problemphyxmini0935-inducedcurrentrunsalongdrawnarrow}
  \lean{PhyXMiniProblems.ProblemPhyXMini0935.InducedCurrentRunsAlongDrawnArrow}
  \uses{def:physics:phyx-mini-0935:phyxminiproblems-problemphyxmini0935-electriccurrentinamperes, def:physics:phyx-mini-0935:phyxminiproblems-problemphyxmini0935-inductionloopsetup}
  Positive signed current means motion along the purple reference arrow
\end{definition}
\begin{proof}
  This is the declared model definition of Induced Current Runs Along Drawn Arrow.
\end{proof}
\begin{definition}[Induced Current Is Clockwise Viewed From Above]
  \label{def:physics:phyx-mini-0935:phyxminiproblems-problemphyxmini0935-inducedcurrentisclockwiseviewedfromabove}
  \lean{PhyXMiniProblems.ProblemPhyXMini0935.InducedCurrentIsClockwiseViewedFromAbove}
  \uses{def:physics:phyx-mini-0935:phyxminiproblems-problemphyxmini0935-traversalsensefromabove, def:physics:phyx-mini-0935:phyxminiproblems-problemphyxmini0935-inductionloopsetup, def:physics:phyx-mini-0935:phyxminiproblems-problemphyxmini0935-inducedcurrentrunsalongdrawnarrow}
  The induced current follows the clockwise traversal seen from above
\end{definition}
\begin{proof}
  This is the declared model definition of Induced Current Is Clockwise Viewed From Above.
\end{proof}
% --- Archon declaration topology end ---
% --- Archon physics formalization source end ---
